
\documentclass[a4paper,10pt]{report}

%\usepackage[landscape]{geometry}
\usepackage[cm-default]{fontspec}
\setromanfont{FreeSerif}
\setsansfont{FreeSans}
\setmonofont{FreeMono}
\usepackage[utf8x]{inputenc}
\usepackage{fontspec}
\usepackage{xunicode}
\usepackage{xltxtra}
\usepackage{xgreek}
\usepackage{amsmath}
\usepackage{unicode-math}
\usepackage{ulem}
\usepackage{color}
\usepackage{verbatim}
\usepackage{nopageno}
\usepackage{graphicx}
\usepackage{textpos}
\setlength{\TPHorizModule}{1cm}
\setlength{\TPVertModule}{1cm}
%\usepackage[colorgrid,texcoord]{eso-pic}
\usepackage[outline]{contour}
\usepackage{wrapfig}
\usepackage{url}
\usepackage{color}
\usepackage{tikz}
\usepackage{fancybox,fancyhdr}
\usepackage{subfigure}
\usepackage{pstricks}
\usepackage{epsfig}
\usepackage{multicol}
\usepackage{listings}
\usepackage{enumerate}
\usepackage{hyperref}
\hypersetup{
  bookmarks=true,
  bookmarksopen=true,
  pdfborder=false,
  pdfpagemode=UseNone,
  raiselinks=true,
  pdfhighlight={/P},
  colorlinks,
  citecolor=black,
  filecolor=black,
  linkcolor=black,
  urlcolor=black
}


\usepackage{fontspec}
\usepackage{xunicode}
\usepackage{xltxtra}


% Margins
%\setlength{\textwidth}{24cm}

%\setlength{\voffset}{0in}
%\setlength{\textheight}{6.8in}
% Colors
\definecolor{cbrown}{rgb}{0.49,0.24,0.07}
\definecolor{cmpez}{rgb}{0.92,0.88,0.79}
\definecolor{cmpez2}{rgb}{0.62,0.33,0.34}
\definecolor{cwhite}{rgb}{1,1,1}
\definecolor{cblack}{rgb}{0,0,0}
\definecolor{cred}{rgb}{0.9,0.15,0.15}
\definecolor{clblue}{rgb}{0.57,0.85,0.97}
\definecolor{corange}{rgb}{0.86,0.49,0.18}
\definecolor{clorange}{rgb}{0.95,0.84,0.61}
\definecolor{cyellow}{rgb}{1,0.95,0.16}
\definecolor{cgreen}{rgb}{0.69,0.8,0.31}
\definecolor{clbrown}{rgb}{0.78,0.67,0.43}
\definecolor{cmagenta}{rgb}{0.79,0.05,0.54}
\definecolor{cgray}{rgb}{0.76,0.73,0.63}
\definecolor{bluesite}{rgb}{0.05,0.43,0.67}
\definecolor{redsite}{rgb}{0.90,0.15,0.15}
\definecolor{backsite}{rgb}{0.07,0.07,0.07}
% ----------- from Costas -----------------------------

%---------------   Main Fonts    -----------------%   

% \setmainfont[Mapping=tex-text]{Calibri}
% \setmainfont[Mapping=tex-text]{Times New Roman}
%\setmainfont[Mapping=tex-text]{Droid Serif}
%\setmainfont[Mapping=tex-text]{Cambria}
%\setmainfont[Mapping=tex-text]{cm-unicode}
% \setmainfont[Mapping=tex-text]{Gentium}
% \setmainfont[Mapping=tex-text]{GFS Didot}
% \setmainfont[Mapping=tex-text]{Comic Sans MS}
 \setmainfont[Mapping=tex-text]{Ubuntu}
% \setmainfont[Mapping=tex-text]{Myriad Pro}
% \setmainfont[Mapping=tex-text]{CMU Concrete}
%\setmainfont[Mapping=tex-text]{DejaVu Sans}
%\setmainfont[Mapping=tex-text]{KerkisSans}
%\setmainfont[Mapping=tex-text]{KerkisCaligraphic}
% \setmainfont[Mapping=tex-text]{Segoe Print}
% \setmainfont[Mapping=tex-text]{Gabriola}

%------------------------------------------------%

% ------------ Mathematics Fonts ----------------%
\setmathfont{Asana-Math.ttf}
% -----------------------------------------------%
% Misc
\author{Κασωτάκη Ε. - Σμαραγδάκης Κ.}
\title{www.math24.gr}
% ---------------- Formation --------------------%


\setlength\topmargin{-0.7cm}
\addtolength\topmargin{-\headheight}
\addtolength\topmargin{-\headsep}
\setlength\textheight{26cm}
\setlength\oddsidemargin{-0.54cm}
\setlength\evensidemargin{-0.54cm}
%\setlength\marginparwidth{1.5in}
\setlength\textwidth{17cm}

\RequirePackage[avantgarde]{quotchap}
\renewcommand\chapterheadstartvskip{\vspace*{0\baselineskip}}
\RequirePackage[calcwidth]{titlesec}
\titleformat{\section}[hang]{\bfseries}
{\Large\thesection}{12pt}{\Large}[{\titlerule[0.9pt]}]
%--------------------------------------------------%
\usepackage{draftwatermark}
\SetWatermarkText{www.math24.gr}
\SetWatermarkLightness{0.9}
\SetWatermarkScale{0.6}
%--------------------------------------------------%


\usepackage{xcolor}
\usepackage{amsthm}
\usepackage{framed}
\usepackage{parskip}

\colorlet{shadecolor}{bluesite!20}

\newtheorem*{orismos}{Ορισμός}

\newenvironment{notation}
  {\begin{shaded}\begin{theorem}}
  {\end{theorem}\end{shaded}}




% Document begins
\begin{document}
\pagestyle{fancy}
\fancyhead{}
\fancyfoot{}
\renewcommand{\headrulewidth}{0pt}
\renewcommand{\footrulewidth}{0pt}


\fancyhead[LO,LE]{
 %\textblockcolor{backsite}
 %\begin{textblock}{5}(-2,-0.55)
  %\rule{0cm}{1cm}
 %\end{textblock}
 \textblockcolor{bluesite}
 \begin{textblock}{5}(-1.5,-0.55)
  \rule{0cm}{1cm}
 \end{textblock}
 %\textblockcolor{bluesite}
 %\begin{textblock}{14}(5.5,-0.55)
 % \rule{0cm}{1cm}
 %\end{textblock}
 \begin{textblock}{0}(-1,-0.25)
 \color{cwhite} \begin{Large}www.math24.gr\end{Large}
 \end{textblock}
\begin{textblock}{0}(16,-0.4)
 \color{cwhite} \includegraphics[height=1.5cm]{math24_logo.png}
 \end{textblock}
\textblockcolor{white}
\begin{textblock}{17}(-1,28)
 \color{backsite} \begin{small}Copyright \textcopyright 2012, Κασωτάκη Ε.(ikasotaki@gmail.com) - Σμαραγδάκης Κ.(kesmarag@gmail.com)\end{small}
 \end{textblock}
}
 
\begin{shaded}
\begin{center}
\huge \textbf{Φυλλάδιο Ασκήσεων}\\
%Πρόσθεση ρητών αριθμών
\end{center} 
\textbf{Μαθηματικά Γ' Γυμνασίου} \hfill \textbf{Ημερομηνία Παράδοσης : \hspace{2em} }
\subsection*{Ονοματεπώνυμο :\hfill  \hspace{5em}}
\end{shaded}
\vspace{2em}
\begin{itemize}
 \item Ορισμός τετραγωνικής ρίζας
 \item Ιδιότητες ριζών
\end{itemize}


\section*{Θεωρία - Τετραγωνική Ρίζα \hfill \small{}}
\begin{itemize}
 \item \textbf{Τετραγωνική ρίζα}  ενός θετικού αριθμό $x$ $(\sqrt{x})$ είναι o θετικός αριθμός που όταν 
        υψωθεί στο τετράγωνο μας δίνει τον αριθμό $x$.\\
\textbf{π.χ} $\sqrt{16}=4$ γιατί $4^{2}=16$
 \item Η τετραγωνική ρίζα του $0$ ισούται με $0$, δηλαδή $\sqrt{0}=0$
 \item Γενικά: $\sqrt{α^{2}}=|α|$ για κάθε πραγματικό αριθμό $α$\\
\textbf{π.χ} $\sqrt{3^{2}}=|3|$ επειδή $\sqrt{3^{2}}=\sqrt{9}=3$ (γιατί $3^{2}=9$)\\
\textbf{π.χ} $\sqrt{(-5)^{2}}=|5|$ επειδή $\sqrt{(-5)^{2}}=\sqrt{25}=5$ (γιατί $5^{2}=25$)
\end{itemize}


\section*{Θεωρία \hfill \small{}}
\textbf{Αν} $β\ge 0$ \textbf{τότε} $(\sqrt{β})^{2}=β$\\
\textbf{π.χ} $(\sqrt{4})^{2}=4$ γιατί $(\sqrt{4})^{2}=(2)^{2}=4$

\section*{Θεωρία - Ιδιότητες Ριζών \hfill \small{}}
\begin{itemize}
 \item $\sqrt{α} \cdot \sqrt{β} = \sqrt{αβ}$ \\
       \textbf{π.χ} $\sqrt{4} \cdot \sqrt{9} = \sqrt{4\cdot 9}$ 
       γιατί $\sqrt{4} \cdot \sqrt{9} =2\cdot 3=6$ και $\sqrt{4\cdot 9}=\sqrt{36}=6$
 \item $\dfrac{\sqrt{α}}{\sqrt{β}}=\sqrt{\dfrac{α}{β}}$ με $β>0$\\
       \textbf{π.χ} $\dfrac{\sqrt{4}}{\sqrt{9}}=\sqrt{\dfrac{4}{9}}$ 
       γιατί $\dfrac{\sqrt{4}}{\sqrt{9}}= \dfrac{2}{3} $ και $\sqrt{\dfrac{4}{9}}=\dfrac{2}{3}$
\end{itemize}













\section*{Άσκηση 1  \hfill \small{30 μονάδες}}
Επιλέξτε τη σωστή απάντηση:
\begin{enumerate}[1)]
 \item το $\sqrt{25}$ ισούται με %1
\begin{multicols}{3}
\begin{enumerate}[i)]
 \item $5$
 \item $-5$
 \item δεν ορίζεται
\end{enumerate}
\end{multicols}
%--------------------------------------
 \item το  $\sqrt{4^{2}}$ ισούται με %t2
\begin{multicols}{3}
 \begin{enumerate}[i)]
 \item $4$
 \item $-4$
 \item δεν ορίζεται
 \end{enumerate}
\end{multicols}
 %-----------------------------------
  \item το $\sqrt{-3^{2}}$ ισούται με %3
\begin{multicols}{3}
\begin{enumerate}[i)]
 \item $3$
 \item $-3$
 \item δεν ορίζεται 
\end{enumerate}
\end{multicols}
%------------------------------
 \item το $\sqrt{(-6)^{2}}$ ισούται με %4
\begin{multicols}{3}
\begin{enumerate}[i)]
 \item $-|-6|$
 \item $|-6|$
 \item δεν ορίζεται
\end{enumerate}
\end{multicols}
%------------------------------------
 \item το $\sqrt{7^{2}}$ ισούται με %5
\begin{multicols}{3}
\begin{enumerate}[i)]
 \item $-|7|$
 \item $|7|$
 \item δεν ορίζεται
\end{enumerate}
\end{multicols}
%---------------------------------------
 \item το $\sqrt{(-8)^{2}}$ ισούται με %6
\begin{multicols}{3}
\begin{enumerate}[i)]
 \item $-|8|$
 \item $|8|$
 \item δεν ορίζεται
\end{enumerate}
\end{multicols}
%-------------------------------------------------
 \item το $\sqrt{3\cdot 4}$ ισούται με %7
\begin{multicols}{3}
\begin{enumerate}[i)]
 \item $4\sqrt{3}$
 \item $3\sqrt{3}$
 \item $\sqrt{3}\cdot \sqrt{4}$
\end{enumerate}
\end{multicols}
%--------------------------------------------
 \item το $\sqrt{4\cdot8}$ ισούται με %8
\begin{multicols}{3}
\begin{enumerate}[i)]
 \item $2\sqrt{8}$
 \item $8\sqrt{4}$
 \item $\sqrt{4}+\sqrt{8}$
\end{enumerate}
\end{multicols}
%--------------------------------------------
 \item το $\sqrt{9\cdot 16}$ ισούται με %9
\begin{multicols}{3}
\begin{enumerate}[i)]
 \item $16\sqrt{9}$
 \item $9\sqrt{16}$
 \item $3\cdot4$
\end{enumerate}
\end{multicols}
%--------------------------------------------
 \item το $\sqrt{5}\cdot \sqrt{6}$ ισούται με %10
\begin{multicols}{3}
\begin{enumerate}[i)]
 \item $\sqrt{5\cdot 6}$
 \item $5\sqrt{6}$
 \item $6\sqrt{5}$
\end{enumerate}
\end{multicols}
%--------------------------------------------
 \item το $\sqrt{2}\cdot \sqrt{7}$ ισούται με %11
\begin{multicols}{3}
\begin{enumerate}[i)]
 \item $7\sqrt{2}$
 \item $2\sqrt{7}$
 \item $\sqrt{14}$
\end{enumerate}
\end{multicols}
%--------------------------------------------
 \item το $\sqrt{2}\cdot \sqrt{50}$ ισούται με %12
\begin{multicols}{3}
\begin{enumerate}[i)]
 \item $50$
 \item $2$
 \item $10$
\end{enumerate}
\end{multicols}
%--------------------------------------------
 \item το $\sqrt{\dfrac{6}{7}}$ ισούται με %13
\begin{multicols}{3}
\begin{enumerate}[i)]
 \item $\dfrac{\sqrt{6}}{\sqrt{7}}$
 \item $\dfrac{\sqrt{6}}{7}$
 \item $\dfrac{6}{\sqrt{7}}$
\end{enumerate}
\end{multicols}
%--------------------------------------------
 \item το $\dfrac{\sqrt{2}}{\sqrt{8}}$ ισούται με %14
\begin{multicols}{3}
\begin{enumerate}[i)]
 \item $\sqrt{\dfrac{2}{8}}$
 \item $\dfrac{\sqrt{2}}{8}$
 \item $\dfrac{2}{\sqrt{8}}$
\end{enumerate}
\end{multicols}
%--------------------------------------------
 \item το $\dfrac{\sqrt{250}}{\sqrt{10}}$ ισούται με %15
\begin{multicols}{3}
\begin{enumerate}[i)]
 \item $5$
 \item $25$
 \item $10$
\end{enumerate}
\end{multicols}
%--------------------------------------------
\end{enumerate}


\section*{Άσκηση 2  \hfill \small{}25  μονάδες}
Να αποδείξετε τις παρακάτω σχέσεις:
\begin{enumerate}[i)]
 \item $\sqrt{40}=2\sqrt{10}$
 \item $\sqrt{125}=5\sqrt{5}$
 \item $\sqrt{32}=4\sqrt{2}$
 \item $\sqrt{3}\cdot \sqrt{12}=6$
 \item $\sqrt{32}\cdot \sqrt{2}=8$
\end{enumerate}



\section*{Άσκηση 3  \hfill \small{}20  μονάδες}
Να υπολογίσετε τις παρακάτω παραστάσεις:
\begin{enumerate}[i)]
 \item $2\sqrt{7}+3\sqrt{7}+6\sqrt{7}$
 \item $5\sqrt{11}-2\sqrt{11}+6\sqrt{11}$
 \item $3\sqrt{3}+5\sqrt{3}-\sqrt{12}$
 \item $4\sqrt{5}-2\sqrt{3}-6\sqrt{5}+8\sqrt{3}$
\end{enumerate}




\section*{Άσκηση 4  \hfill \small{} 25 μονάδες}
Να μετατρέψετε τα παρακάτω κλάσματα που έχουν άρρητους παρονομαστές σε ισοδύναμα κλάσματα με ρητούς παρονομαστές:
\begin{enumerate}[i)]
 \item $\dfrac{3}{\sqrt{5}}$
 \item $\dfrac{1}{\sqrt{7}}$
 \item $\dfrac{\sqrt{12}}{\sqrt{3}}$
 \item $\dfrac{\sqrt{2}}{\sqrt{8}}$
 \item $\dfrac{\sqrt{3}}{\sqrt{7}}$
\end{enumerate}



 









%\let\thefootnote\relax\footnotetext{Επιμέλεια : Κασωτάκη Ειρήνη (ikasotaki@gmail.com),    Σμαραγδάκης Κώστας (kesmarag@gmail.com)} 


\end{document}
